% ---
% Reemplaza líneas 1618-1668 de Docs/Technical_Report.tex
% Decisiones: D-014, H-011, H-054
% ---

\subsubsection{Integración cuántico-clásica: \textit{embeddings} precomputados}

La integración entre los circuitos cuánticos y el clasificador clásico se formula
mediante una arquitectura desacoplada basada en representaciones intermedias.
En lugar de ensamblar el circuito cuántico como una
capa diferenciable dentro de la red neuronal mediante adaptadores de software como
\texttt{TorchConnector} \cite{qiskitml2024}, el pipeline separa la fase
de transformación cuántica de la fase de clasificación supervisada.

El flujo de procesamiento opera de acuerdo con las siguientes etapas metodológicas:
\begin{enumerate}
  \item \textbf{Generación determinista de representaciones:} A partir del vector de
        características seleccionado y escalado $\mathbf{x} \in [0, \pi]^{12}$ (M4),
        el circuito cuántico compuesto por el mapa de características (\texttt{zz\_feature\_map}
        en C4 o \texttt{pauli\_feature\_map} en C5) y el \textit{ansatz} \texttt{real\_amplitudes}
        con parámetros fijos $\boldsymbol{\theta}$ calcula de forma analítica el vector
        de expectativas exactas $\mathbf{e}(\mathbf{x}, \boldsymbol{\theta}) \in [-1, 1]^{12}$.
  \item \textbf{Almacenamiento de representaciones:} Los vectores resultantes se
        precomputan una sola vez para la totalidad del conjunto de datos y se persisten
        en disco en formato estructurado, eliminando la necesidad de reevaluar circuitos
        durante el ciclo de optimización del clasificador.
  \item \textbf{Estandarización:} Los vectores de cada condición experimental
        se estandarizan empleando la media y desviación estándar ajustadas sobre el conjunto de entrenamiento de cada partición, compensando la disparidad
        en las escalas numéricas de las distintas condiciones (M7).
  \item \textbf{Entrenamiento del clasificador clásico:} El perceptrón multicapa (MLP)
        en PyTorch se entrena sobre las representaciones estandarizadas,
        optimizando sus pesos sin modificar los parámetros del circuito cuántico.
\end{enumerate}

Este diseño responde a tres razones:
\begin{itemize}
  \item \textbf{Validez interna del experimento:} Las condiciones clásicas C2 (PCA) y
        C3 (Kernel PCA) operan como transformaciones no supervisadas que no acceden a las
        etiquetas de patología para construir el nuevo espacio. Si C4 y C5 ajustaran
        $\boldsymbol{\theta}$ contra las etiquetas durante el entrenamiento, cualquier
        ganancia en clasificación podría atribuirse al ajuste paramétrico supervisado y no
        a las propiedades geométricas de la codificación cuántica $\phi$.
  \item \textbf{Correspondencia con la pregunta de investigación:} El objetivo del trabajo
        consiste en determinar en qué medida el mapeo cuántico al espacio de Hilbert modifica
        la separabilidad intrínseca de las características radiómicas, aislando la capacidad
        representacional del \textit{feature map} frente a la optimización variacional.
  \item \textbf{Viabilidad computacional:} Evita el costo de cientos de horas de ejecución
        que demandaría el método \textit{parameter-shift} en cada época de entrenamiento; en el simulador la retropropagación sobre el vector de estado abarataría el entrenamiento, pero ese atajo no existe en hardware (sección~\ref{sec:parameter_shift}).
\end{itemize}
